\documentclass[12pt]{article}
\usepackage[utf8]{inputenc}
\usepackage[T1]{fontenc}
\usepackage{hyperref}
\title{Labo Phowa Language Resources for K-12 Curriculum Development}
\author{Compiled by Claude Sonnet 5.5 (claude-sonnet-5-5) for Andr\'as Kornai}
\date{}
\begin{document}
\maketitle

\section{Introduction}
Labo Phowa (ISO 639-3 ypb, Glottocode labo1243) is a Loloish (Sino-Tibetan) language of southern China, usually grouped with the Phula varieties. Wikipedia treats Phowa as a dialect cluster with three main varieties (Hlepho, Ani and Labo) that may be distinct languages, and gives about 17,000 Labo speakers (2011). Joshua Project and GRN give ethnic-group populations of about 21,000 and 22,000. Sources disagree on these figures. Glottolog lists the code as the result of a 2008 split of the earlier ``Puwa Yi'' entry. Very little has been published specifically on Labo. This dossier lists only what the searches returned.

\section{Primary Educational Resources}
No primers, textbooks or school programmes for Labo Phowa turned up in the searches. The language is not known to be taught in schools.

\section{Linguistic and Reference Materials}
\begin{itemize}
\item Phowa language (Wikipedia overview of the Hlepho, Ani and Labo varieties): \url{https://en.wikipedia.org/wiki/Phowa_language}
\item Pelkey, ``Puzzling over Phula'' (academic survey of the Phula varieties; it notes how little has been documented): \url{https://www.academia.edu/2435607/Puzzling_over_Phula}
\item Glottolog entry for ypb: \url{https://glottolog.org/resource/languoid/iso/ypb}
\item Glottolog entry for the earlier Puwa Yi grouping: \url{https://glottolog.org/resource/languoid/id/puwa1234}
\item Wikipedia on Southeastern Loloish languages (classification context): \url{https://www.Wikipedia.com/wiki/Phula_language}
\end{itemize}

\section{Religious and Community Resources}
\begin{itemize}
\item Joshua Project profile of Labo Phowa in China (Joshua Project reports no Bible translation): \url{https://joshuaproject.net/people_groups/20609}
\end{itemize}

\section{Digital and Audio Resources}
\begin{itemize}
\item Global Recordings Network language page for ypb: \url{https://globalrecordings.net/en/language/ypb}
\end{itemize}
Only the GRN page was found. I did not verify what recordings it holds.

\section{Curriculum Development Context}
Searches found no orthography, dictionary, grammar or teaching materials for Labo Phowa. A native-language curriculum would first need community-led documentation. The academic literature on Phula is the best starting point. Mandarin is the language of schooling.
\end{document}